\begin{lemma}[Step 1, Subcase 2.2: the third end is inside the starred predecessor]
Let $G,e_0,[k],T,\varphi$ be a finite loopless multigraph with exact edge
identities, distinguished uncolored edge, color set, positive-rung
certified ETT, and generating coloring.  Assume maximum property,
properness, exact one-uncoloredness, and the complete six-part lower-rung
package, whose assertion (iii) gives unconditional stability before each PE
support, predecessor-ladder stability only when its first-appearance index is
greater than one, and unconditional equal-index stored-pair disjointness.  Let $H$ be a good hierarchy over a hierarchy base, with
$d<q_0$ dividers.  Assume the same-rung fewer-than-$q_0$ elementarity
hypothesis for arbitrary certified inputs and generating colorings, the
large bound $|\overline\varphi(T_n)|\geq2n+11$, the separate terminal
certificate
\[
 T\text{ is }
 \left(\bigcup_{\eta\in D^H_{n,d+1}}\Gamma^d_\eta\right)^-
 \text{-closed under }\varphi,
\]
and elementarity of $V(T)$.  Put $A=H.\mathit{starred}(d)$.

Let $\alpha\ne\beta$, with $\alpha$ missing and singleton-closed on $T$
and $\beta$ not missing on $T$.  Let $P_1,P_2,P_3$ be pairwise
support-disjoint exact $(T,\varphi,\{\alpha,\beta\})$-exit paths.  Their
outside ends $a_1,a_2,a_3$ and tree ends $b_1,b_2,b_3$ are pairwise
distinct, with $b_1\prec_T b_2\prec_T b_3$.  Retain the exact
$\beta$-colored boundary-edge identity $f_j$ of each path incident with
$b_j$.  Suppose $b_3\in V(A)$.

Assume that a selected color $\theta\ne\alpha,\beta$ is missing and
singleton-closed on $A$ and singleton-closed on $T$.  Assume also the
explicit no-smaller-presentation consequence of the global minimum: for
every proper, exactly one-uncolored coloring $\sigma$ stable from
$\varphi$ on the final ladder tree, arbitrary replacement input,
corresponding same-rung maximum ETT $T'$, good hierarchy with fewer than
$d$ dividers, its own large estimate, and its separate terminal
minus-closed certificate, there do not exist two support-disjoint exact
exit legs with distinct outside ends for a missing singleton-closed first
color.  Then the displayed Subcase 2.2 data are contradictory.
\end{lemma}
\begin{proof}
Put $d=H.\mathit{dividerCount}$.  If $d>0$, the contradiction is
\noderef{Step1CaseTwoInsidePositiveDivider}.  If $d=0$, decide the exact
terminal extension-kind equation.  The PE alternative is
\noderef{Step1CaseTwoInsideZeroPE}, and its complement is
\noderef{Step1CaseTwoInsideZeroNonPE}.  All three lemmas use the displayed
three edge-identity-preserving exit paths and the same full replacement-family
minimality hypothesis, so the alternatives exhaust the present data.
\end{proof}
